\section{Closing the contradiction and the maximal estimate}
\label{sec:closure}

The stationary packets require the limiting adjusted rate
\(\mathcal M-w(g-d)\) to be at least \(\tau+B_*\).
We show that the projection constraints are incompatible with this demand
and the bounds for \(g\) on thin offset stripes.  The strong compactness of
\(v_1\) allows us to use the constraints on the sets where its limiting
value is positive.  This excludes a positive critical exponent.
We then pass from the critical inequality to the maximal estimate,
retaining the shading-density power through the reduction to arbitrary
\(L^3\) functions.

\subsection{Excluding the stationary rates}

\begin{proposition}[Vanishing at the maximal endpoint]
\label{prop:critical-vanishing}
There are values \(p>2\) arbitrarily close to \(2\) for which
\(h(p,0)=0\).
\end{proposition}

\begin{proof}
Suppose otherwise.  Choose an open interval \(I\) of \(p\)-values, with
lower endpoint \(p_{\min}>2\), on which \(h(p,0)>0\).  Fix the smoothing
parameter before making any of the subsequent selections, with
\begin{equation}
  0<\eta\le
  \min\left\{\frac1{32},\frac{p_{\min}-2}{16}\right\}.
  \label{eq:closure-eta}
\end{equation}
This choice is legitimate because the critical inequality at \(q=0\) is
independent of \(\eta\).  Lemma~\ref{lem:critical-point} supplies an interior
differentiability point with \(h>0\) and \(h_z>0\).  The construction in
Proposition~\ref{prop:canonical}, Proposition~\ref{prop:stationary-data}, and
Lemma~\ref{lem:angular-frostman} then supplies fixed parameters satisfying
\begin{gather}
  p>2,\qquad 0\le q<1,\qquad d=2-q,\qquad
  0<h\le1-q,\qquad x=1-q-h<1, \label{eq:closure-parameters-a}\\
  0<\tau<1,\qquad \ell=1-\tau,\qquad
  0<\beta<1,\qquad s_0=1-\beta, \label{eq:closure-parameters-b}\\
  B_*=h+\ell\bigl(p\beta-\mathcal D(q,\beta)\bigr),\qquad
  w=h_z+\ell\mathcal D_q(q,\beta),\qquad 0<w\le1.
  \label{eq:closure-parameters-c}
\end{gather}
In particular, Lemma~\ref{lem:angular-frostman} supplies the strict
inequality \(s_0>0\).

Recall that $g$ is the total matrix entropy rate and $\mathcal M$ its
speed-weighted rate from \eqref{eq:entropy-group-rates}.
For brevity in this proof put
\begin{equation}
  b_+=(\beta-q)_+,\qquad T_*=\tau+B_*,\qquad
  \mathcal Q=\mathcal M-w(g-d),
  \label{eq:closure-demand-notation}
\end{equation}
where \(r_+=\max\{r,0\}\).  The smoothing definition in \eqref{eq:temporal-integrand} gives
\begin{equation}
  0\le\mathcal D(q,\beta)\le\min\{q,\beta/2\}.
  \label{eq:closure-cost-bound}
\end{equation}
If \(q<(1/2-\eta)\beta\), then the integrand defining \(\mathcal D\)
equals \(q\) throughout its averaging interval.  Consequently
\begin{equation}
  \mathcal D(q,\beta)=q,\qquad
  \mathcal D_q(q,\beta)=1,\qquad w>\ell.
  \label{eq:closure-low-q}
\end{equation}

We record precisely the limiting statements to be used.  Fix a compact
interval \(J\subset(0,1)\).  Proposition~\ref{prop:entropy-demand} and
Lemma~\ref{lem:stripe-demand} imply that every weak-star limit of
\(\mathcal Q(\cdot,a)\), along generic offsets \(a\to0\), is at least
\(T_*\) almost everywhere on \(J\).  The same lower bound holds for every
weak-star limit of stripe averages.  Moreover, the right stripe averages
of \(g\) are at least \(d\), and the left stripe averages are at most
\(d\), whenever the latter offsets are available.  In formulas,
\begin{equation}
  \frac1\epsilon\int_0^\epsilon g(c,a)\,da\ge d,
  \qquad
  \frac1\epsilon\int_{-\epsilon}^0 g(c,a)\,da\le d
  \label{eq:closure-stripes}
\end{equation}
for almost every \(c\in J\) and sufficiently small \(\epsilon>0\), on
the sides being used.  These inequalities may be integrated against every
nonnegative \(L^1\) test function supported on \(J\).  The rates are
uniformly bounded, so weak-limit inequalities initially established with
smooth tests extend to such tests by \(L^1\) approximation.  We will use
this extension to localize on measurable subsets of \(J\).

For \(w<1\), Lemma~\ref{lem:velocity-trace} gives strong one-sided
\(L^1(J)\) traces for \(v_1\) through a fixed full-measure set of generic
offsets.  The corresponding stripe averages converge to the same trace.
For \(w=1\), that lemma gives strong subsequential \(L^1(J)\)
compactness along every sequence of generic offsets.  We use this strong
convergence to localize the constraints that hold where \(v_1>0\);
weak-star convergence suffices for the other rates.

By Corollary~\ref{cor:normal-position-rate}, on the right we also have
\begin{equation}
  a_1\le x=1-q-h<1.
  \label{eq:closure-x-rate}
\end{equation}
It follows in particular that
\begin{equation}
  (a_1-s_0)_+\le(\beta-q-h)_+\le b_+.
  \label{eq:closure-extra-rate}
\end{equation}

The domain \(\mathcal U\) in \eqref{eq:entropy-domain} contains small
offsets of both signs uniformly on \(J\) when \(w<1\), and small
positive offsets when \(w=1\).  Thus
Proposition~\ref{prop:projection-rates} applies on every stripe used
below.  The bound \eqref{eq:closure-x-rate} is available on the right.

The choice of stripe follows from the identity
\[
 \mathcal Q=(\mathcal M-rg)+rd+(r-w)(g-d),
 \qquad r\in\R.
\]
After averaging, an upper bound for \(\mathcal M-rg\) gives an upper bound
for \(\mathcal Q\) on the right when \(r\le w\), and on the left
when \(r\ge w\).  For distinct speeds the projection inequalities
pair \(a_1\) with \(y_1\) and \(u_1\) with \(v_1\), making
\(r=\tau/2\) useful.  Direct bounds for \(\mathcal M\), corresponding
to \(r=0\), handle two further regimes.  Table~\ref{tab:closure-cases}
records the five cases.  In the last case, compactness first gives a
gap common to all weak limits of generic slices; averaging then gives
the stripe contradiction.

\begin{table}[htbp]
\centering
\small
\begin{tabular}{@{}llll@{}}
\toprule
Speeds & Parameter regime & Stripe & Main comparison\\
\midrule
Distinct & $w\ge\tau/2$ & Right & $\mathcal M-\tau g/2$\\
Distinct & $w<\tau/2$, $\tau<2\ell$ & Right & $\mathcal M$\\
Distinct & $w<\tau/2$, $\tau\ge2\ell$ & Left
 & $\mathcal M-\tau g/2$, trace of $v_1$\\
Tied & $q\ge(1/2-\eta)\beta$ & Right & $\mathcal M$\\
Tied & $q<(1/2-\eta)\beta$ & Right
 & $\mathcal M-g/2$, compactness of $v_1$\\
\bottomrule
\end{tabular}
\caption{The comparisons used to exclude the stationary rates.
The stripe bounds for $g$ enter only after averaging.}
\label{tab:closure-cases}
\end{table}

\paragraph{Distinct speeds, \(w\ge\tau/2\).}
Suppose first that \(\tau\ne\ell\).  In coordinate notation,
\[
  g=a_1+y_1+u_1+v_1,
  \qquad \mathcal M=\tau y_1+\ell u_1+v_1.
\]
On the right, Proposition~\ref{prop:projection-rates} and
\eqref{eq:closure-extra-rate} give
\(a_1\ge y_1\), \(u_1\ge v_1\), and \(u_1,v_1\le b_+\).  Thus
\begin{align}
  \mathcal M-\frac\tau2g
   &=-\frac\tau2(a_1-y_1)
      +\left(\ell-\frac\tau2\right)u_1
      +\left(1-\frac\tau2\right)v_1
      \le2\ell b_+.
      \label{eq:closure-first-pointwise}
\end{align}
If \(\ell-\tau/2\ge0\), bound both last rates by \(b_+\).  If
\(\ell-\tau/2<0\), first use \(u_1\ge v_1\); the resulting coefficient
of \(v_1\) is \(\ell+1-\tau=2\ell>0\).  This verifies
\eqref{eq:closure-first-pointwise} for either ordering of the middle speeds.
Since \(\tau/2-w\le0\), the right stripe inequality in
\eqref{eq:closure-stripes} implies that every upper stripe limit of
\(\mathcal Q\) is at most
\[
  U_1=\frac\tau2d+2\ell b_+.
\]
Its gap below the required lower limit is
\[
  T_*-U_1
   =h+\frac{\tau q}{2}
      +\ell\bigl(p\beta-\mathcal D(q,\beta)-2b_+\bigr).
\]
For \(q<\beta\), this is at least
\(h+\tau q/2+\ell((p-2)\beta+q)>0\).  For \(q\ge\beta\), it is at
least \(h+\tau q/2+\ell(p-1/2)\beta>0\).  Either possibility
contradicts the lower stripe limit \(T_*\).

\paragraph{Distinct speeds, \(w<\tau/2\) and \(\tau<2\ell\).}
Here \(\ell>1/3\).  If \(q<(1/2-\eta)\beta\),
\eqref{eq:closure-low-q} would give \(w>\ell>\tau/2\).  We must
therefore have \(q\ge(1/2-\eta)\beta\).  On the right,
\(y_1\le a_1\le1-q\) and \(u_1,v_1\le b_+\), so
\[
  \mathcal M\le U_2:=\tau(1-q)+(1+\ell)b_+.
\]
The right stripe bound and \(w>0\) make \(U_2\) an upper stripe
limit for \(\mathcal Q\) as well.  If \(q<\beta\), subtraction gives
\begin{align*}
  T_*-U_2
    &=h+\bigl(\ell(p-1)-1\bigr)\beta
           -\ell\mathcal D(q,\beta)+2q\\
    &\ge h+\bigl(\ell(p-3/2)-2\eta\bigr)\beta>0.
\end{align*}
The last inequality follows from \(\ell>1/3\), \(p>2\), and
\(\eta\le1/32\); its coefficient is greater than
\(1/6-1/16>0\).  If \(q\ge\beta\), the gap is at least
\(h+\tau q+\ell(p-1/2)\beta>0\).  This again contradicts the demand.

\paragraph{Distinct speeds, \(w<\tau/2\) and \(\tau\ge2\ell\).}
Use the left stripe.  Since \(w<\tau/2<1/2\), its offsets are
available and \(v_1\) has a strong left trace, denoted by \(v_*\).
Take a weak-star limit \(\bar g\) of left stripe averages.  The
projection constraints imply, wherever \(v_1>0\),
\begin{equation}
  g\ge2\min\{1,s_0+v_1\}+2v_1.
  \label{eq:closure-positive-v-lower}
\end{equation}
We give the localization estimate, since convergence of the averaged rate
alone would not suffice here.  The full generic-slice trace implies
\begin{equation}
 R_\epsilon:=\frac1\epsilon\int_{-\epsilon}^0
       \|v_1(\cdot,a)-v_*\|_{L^1(J)}\,da\longrightarrow0.
 \label{eq:closure-mixed-trace-error}
\end{equation}
Indeed the integrand tends to zero through a fixed full-measure set of
offsets, so its essential supremum over \((-\epsilon,0)\) tends to zero.
For \(E_\delta=\{c\in J:v_*(c)\ge\delta\}\), where \(\delta>0\), let
\(B_a=E_\delta\cap\{c:v_1(c,a)=0\}\).  Then
\begin{equation}
 |B_a|\le\delta^{-1}\|v_1(\cdot,a)-v_*\|_{L^1(J)}.
 \label{eq:closure-bad-zero-set}
\end{equation}
The function \(f_0(t)=2\min\{1,s_0+t\}+2t\), for \(0\le t\le1\),
is bounded by \(4\) and has Lipschitz constant at most \(4\).
For every bounded nonnegative test \(\varphi\) supported on \(E_\delta\),
\eqref{eq:closure-positive-v-lower} and \(g\ge0\) consequently give
\[
 \int_J\varphi(c)\frac1\epsilon\int_{-\epsilon}^0g(c,a)\,da\,dc
 \ge \int_J\varphi(c)f_0(v_*(c))\,dc
       -4\|\varphi\|_\infty(1+\delta^{-1})R_\epsilon.
\]
Take the stripe weak-star limit, then let \(\delta=1/m\downarrow0\).
The resulting localized inequality is
\[
  \bar g\ge2\min\{1,s_0+v_*\}+2v_*
  \quad\hbox{almost everywhere on }\{v_*>0\}.
\]
On the other hand, \(\bar g\le d\le2\).  If \(s_0+v_*\ge1\), the
displayed lower bound is greater than \(2\).  Otherwise it is
\(2-2\beta+4v_*\).  It follows that
\begin{equation}
  v_*\le\beta/2\quad\hbox{almost everywhere},
  \label{eq:closure-left-v-trace}
\end{equation}
including the set where the trace vanishes.

Because \(\ell-\tau/2\le0\), the inequalities \(a_1\ge y_1\) and
\(u_1\ge v_1\) also imply
\[
  \mathcal M-\frac\tau2g\le2\ell v_1.
\]
Now \(\tau/2-w>0\).  Combining the left stripe bound for \(g\) with
\eqref{eq:closure-left-v-trace}, every upper stripe limit of
\(\mathcal Q\) is bounded by
\[
  U_3=\frac\tau2d+\ell\beta.
\]
But
\[
  T_*-U_3
   =h+\frac{\tau q}{2}
       +\ell\bigl((p-1)\beta-\mathcal D(q,\beta)\bigr)
   \ge h+\frac{\tau q}{2}+\ell(p-3/2)\beta>0.
\]
This excludes the last distinct-speed case, including \(\tau=2\ell\).

\paragraph{Tied speeds with \(q\ge(1/2-\eta)\beta\).}
It remains to consider \(\tau=\ell=1/2\).  Write \(m_1\) for the
middle-group rate.  Then
\[
  g=a_1+m_1+v_1,\qquad \mathcal M=m_1/2+v_1.
\]
On the right, Proposition~\ref{prop:projection-rates} yields
\(m_1\le1+b_+\) and \(v_1\le b_+\).  Therefore
\[
  \mathcal M\le U_4:=\frac12+\frac32b_+.
\]
This is again an upper stripe limit for \(\mathcal Q\).  For
\(q<\beta\), its gap below demand satisfies
\[
  T_*-U_4
    \ge h+\left(\frac{p-2}{2}-\frac{3\eta}{2}\right)\beta>0.
\]
Indeed \eqref{eq:closure-eta} bounds the coefficient below by
\(13(p-2)/32\).  For \(q\ge\beta\), the gap is at least
\(h+(p-1/2)\beta/2>0\).  This case is impossible as well.  In
particular, the same \(\eta\) fixed before the canonical construction
works independently of the subsequently selected \(h,\tau,\ell,\beta\).

\paragraph{Tied speeds with \(q<(1/2-\eta)\beta\).}
Here \(\mathcal D(q,\beta)=q\) and \(w>1/2\).  Put
\begin{equation}
  U_0=\frac12+\frac{\beta-q}{2},\qquad
  K_0=T_*-U_0=h+\frac{p-1}{2}\beta>0.
  \label{eq:closure-tied-gap}
\end{equation}
Take any sequence of generic positive offsets tending to zero.  On a
fixed compact interior interval \(J\), extract a subsequence for which
\(v_1\to v_*\) strongly in \(L^1(J)\) and the other bounded rates
converge weak-star.  This is allowed also when \(w=1\).  Denote their
limits by bars.  The demand gives \(\overline{\mathcal Q}\ge T_*\).

On \(E_\delta=\{v_*\ge\delta\}\), use
\eqref{eq:closure-bad-zero-set} along the selected sequence of slices:
the zero-set measure is at most
\(\delta^{-1}\|v_{1,n}-v_*\|_1\to0\).  Since \(a_1<1\) on the right, positive \(v_1\)
and the projection constraint \(v_1\le(a_1-s_0)_+\) force
\(a_1\ge s_0+v_1\).  Thus
\[
  \mathcal M-g/2=(v_1-a_1)/2\le-s_0/2
\]
on the part with positive \(v_1\).  Explicitly, for bounded
\(\varphi\ge0\) supported on \(E_\delta\),
\[
 \int_J\varphi(a_{1,n}-v_{1,n}-s_0)
  \ge-s_0\|\varphi\|_\infty\delta^{-1}\|v_{1,n}-v_*\|_1
  \longrightarrow0\quad\hbox{from below}.
\]
Weak-star convergence of \(a_{1,n}\), strong convergence of \(v_{1,n}\),
and then \(\delta=1/m\downarrow0\) justify the inequality on
\(\{v_*>0\}\).  Passing to this localized limit gives
\begin{align*}
  T_*\le\overline{\mathcal Q}
    &\le\frac d2-\frac{s_0}{2}
               +(w-1/2)(d-\bar g)\\
    &=U_0+(w-1/2)(d-\bar g).
\end{align*}
Consequently \(\bar g\le d-K_0/(w-1/2)\) on \(\{v_*>0\}\).
On \(\{v_*=0\}\), the bounds \(m_1\le1+\beta-q\) and strong
convergence of \(v_1\) give \(\overline{\mathcal M}\le U_0\).
The same demand then gives
\(T_*\le U_0+w(d-\bar g)\), and hence
\(\bar g\le d-K_0/w\).  Since \(w-1/2<w\), the two parts imply
the common bound
\begin{equation}
  \bar g\le d-\frac{K_0}{w}
  \quad\hbox{almost everywhere on }J.
  \label{eq:closure-uniform-tied-gap}
\end{equation}
Its positive gap depends only on the fixed parameters, not on the chosen
sequence or the subsequential trace.

We now convert the common gap in
\eqref{eq:closure-uniform-tied-gap}, obtained from arbitrary offset
sequences, to a stripe contradiction.
Fix a nonnegative test \(\varphi\in L^1(J)\) with positive integral.
If arbitrarily small generic positive offsets satisfied
\[
  \int_J\varphi(c)g(c,a)\,dc
    >\left(d-\frac{K_0}{2w}\right)\int_J\varphi(c)\,dc,
\]
choose a sequence of them tending to zero.  The preceding compactness
argument would give a weak-star limit violating
\eqref{eq:closure-uniform-tied-gap}.  Thus all sufficiently small generic
positive offsets satisfy the reverse inequality.  Averaging it over
\(0<a<\epsilon\) contradicts the first inequality of
\eqref{eq:closure-stripes}; the exceptional offsets form a null set.

Every possible choice of the stationary parameters has now been excluded.
This contradicts the differentiability point obtained under the assumed
failure of critical vanishing, and proves the proposition.
\end{proof}

\subsection{From the critical inequality to maximal functions}

The following reduction uses only the indexed model, its regularization
lemma, and the critical inequality at \(q=0\).  In particular its proof
is independent of the stationary construction and of
Proposition~\ref{prop:critical-vanishing}.

\begin{proposition}[Maximal reduction]
\label{prop:maximal-reduction}
Let \(p>2\) be in the parameter range of Definition~\ref{def:critical},
and suppose that \(h(p,0)=0\).  For every \(\kappa>0\) there is a finite
constant \(C_{p,\kappa}\) such that, for all \(0<\delta<1\) and all
complex-valued \(f\in L^3(\R^3)\),
\begin{equation}
  \|K_\delta f\|_{L^3(S^2,\sigma)}
    \le C_{p,\kappa}\,
       \delta^{-(p-2)/3-\kappa}\|f\|_{L^3(\R^3)}.
  \label{eq:closure-quantitative-maximal}
\end{equation}
\end{proposition}

\begin{proof}
We first establish a uniform discrete estimate, then an indicator estimate,
and finally \eqref{eq:closure-quantitative-maximal}.  All geometric constants
in this proof are absolute unless their dependence is displayed.

\paragraph{A discrete estimate in one bounded chart.}
Let \(N\ge2\) be dyadic.  Consider unit-weight lines in a bounded matrix
chart, with pairwise \(c/N\)-separated spatial slopes, where \(c>0\) is
a fixed constant.  Let \(A\) be a finite collection of \(N^{-1}\)-cubes.
Suppose each of \(n\) such lines has at least
\(c_1\max\{1,\lambda N\}\) marked time bins whose centerline points
lie within \(C_1/N\) of a cube in \(A\).  Here \(0<\lambda\le1\)
and \(c_1,C_1\) are fixed constants.  Then, for each \(\xi>0\),
\begin{equation}
  |A|\ge c_{p,\xi}\,N^{-\xi} nN
                    \max\{N^{-1},\lambda\}^{p+1},
  \label{eq:closure-discrete-set}
\end{equation}
where \(|A|\) is a cell count.  Fixed bounded changes of the chart,
of the resolution, and of the neighboring-cell allowance only change
the constant.

Here is the uniformity argument.  For a full-time plank of widths
\(a,b\) in finest cell units, the slopes of lines contained in it lie
in a rotated rectangle with side lengths \(O(a/N)\) and \(O(b/N)\).
This follows by comparing the two endpoint cross-sections of the plank.
Since \(1\le a\le b\le N\), slope separation bounds their number by
\(O(ab)\).  Consequently
\begin{equation}
  \Delta_0=\sup_B\frac{n_B}{ab}\le C
  \label{eq:closure-separated-clustering}
\end{equation}
for the original family and all its subfamilies.

If \eqref{eq:closure-discrete-set} failed uniformly for a fixed
\(\xi>0\), choose a violating sequence with \(N\to\infty\).  Bounded
\(N\) cannot cause this failure, because \(n\le CN^2\) and any
nonempty \(A\) has at least one cell.  On each line retain exactly
\(k\asymp N\theta\) of the allowed bins, where
\(\theta=\max\{N^{-1},\lambda\}\); integer rounding is harmless.
The union of their event cells lies in a fixed enlargement of \(A\),
of cell count at most \(C|A|\).  Along a subsequence write
\(\theta=N^{-\alpha+o(1)}\), with \(0\le\alpha\le1\).
Apply Lemma~\ref{lem:regularization} to this indexed event graph.  It
retains \(nkN^{-o(1)}\) events in weighted count and gives a limiting
profile \(F\).  The number of events per active index is at most \(k\),
whereas retained index mass is at most \(n\).  Total incidence retention
therefore forces this number to be \(kN^{-o(1)}\) in exponent.  It
follows that \(F(1)=\alpha\).  The retained graph has raw multiplicity
at least
\[
  \frac{nkN^{-o(1)}}{C|A|}
      \ge N^{p\alpha+\xi-o(1)}
\]
along a sequence violating \eqref{eq:closure-discrete-set}.
At \(q=0\), the temporal cost of every profile is exactly
\(\Pi_F(1)=pF(1)=p\alpha\).  The critical inequality with
\(h(p,0)=0\) and \eqref{eq:closure-separated-clustering} bounds this
same multiplicity by \(N^{p\alpha+o(1)}\), a contradiction.  This
proves \eqref{eq:closure-discrete-set} with a constant uniform over
all the finite configurations under discussion.

\paragraph{Cylinder geometry and translations.}
Use an auxiliary family of cylinders of length \(4\) and radius
\(4/N\), and let \(\mathcal K_N\) denote the supremum of their
normalized averages over all centers, as a function of direction.
Thus \(0\le\mathcal K_N\one_E\le1\) for every measurable set \(E\).
Let \(E\) be a finite union of cubes of the fixed \(N^{-1}\) grid,
and suppose one of these cylinders has average of \(\one_E\) greater
than \(\lambda>0\).

Cover \(S^2\) by finitely many fixed small direction charts.  After an
axis permutation and, if needed, reversal, the time component in each
chart is bounded away from zero.  A cylinder in such a chart meets only
\(O(1)\) fine spatial cubes in each time bin of length comparable to
\(N^{-1}\).  Its intersection with any one bin has volume
\(O(N^{-3})\), while its full volume is comparable to \(N^{-2}\).
An average greater than \(\lambda\) therefore supplies at least
\(c\lambda N\) occupied time bins, and at least one if this quantity
is smaller than a constant.  The centerline point in each such bin is
within \(C/N\) of an occupied cube of \(E\).  These are the events
needed in \eqref{eq:closure-discrete-set}.

To allow arbitrary translations, group cylinder centers into unit grid
cubes \(Q\).  Every cylinder in a group lies in a fixed enlargement
\(Q^*\), and the regions \(Q^*\) have bounded overlap as \(Q\) varies.
For one group and one direction chart, all centerlines have bounded
intercepts and slopes.  Extend their lines across one fixed common
time interval containing their cylinder traces, then translate and
rescale that interval and the spatial chart by fixed factors to fit the
indexed model.  Only the marked events on the original cylinder are
used.  A comparable dyadic resolution and bounded neighboring-cell
enlargements account for these fixed changes.  Thus
\eqref{eq:closure-discrete-set} applies to the cubes of \(E\) meeting
\(Q^*\), with constants independent of the location of \(Q\).
Pairwise angular separation in a fixed chart is equivalent, up to
fixed constants, to separation of the spatial slopes.

\paragraph{A directional distribution bound.}
For \(0<\lambda<1\), let
\[
  S_\lambda=\{\omega\in S^2:
                         \mathcal K_N\one_E(\omega)>\lambda\}.
\]
In each fixed chart, choose a maximal \(c/N\)-separated subset of
\(S_\lambda\), of cardinality \(n\), and choose for each selected
direction a cylinder witnessing the strict inequality.  Such a finite
maximal subset exists by the packing bound on the sphere.  The
\(C/N\)-caps around its elements cover the high-value directions in
that chart, so their measure is at most \(CnN^{-2}\).

Apply \eqref{eq:closure-discrete-set} to the spatial groups just
described and sum over the groups.  Their ambient regions overlap a
bounded number of times, whence
\[
  N^3|E|
    \ge c_{p,\xi}N^{-\xi}nN
                     \max\{N^{-1},\lambda\}^{p+1},
\]
where \(|E|\) now denotes Lebesgue volume.  Sum over the finitely many
direction charts.  The cap measure bound gives
\begin{equation}
  \sigma(S_\lambda)
    \le C_{p,\xi}N^\xi |E|
                    \max\{N^{-1},\lambda\}^{-(p+1)}.
  \label{eq:closure-indicator-distribution}
\end{equation}
In particular, for \(N^{-1}\le\lambda<1\),
\begin{equation}
  \lambda^3\sigma(S_\lambda)
       \le C_{p,\xi}N^{p-2+\xi}|E|.
  \label{eq:closure-weak-three}
\end{equation}
The factor \(N^{p-2}\) comes from
\(\lambda^{2-p}\le N^{p-2}\).  Thus the density cost has been
retained throughout this reduction.

For a nonempty union of fine grid cubes, \(|E|\ge N^{-3}\).  The
distribution formula, \(\sigma(S^2)=4\pi\), and
\eqref{eq:closure-weak-three} imply
\begin{align}
  \|\mathcal K_N\one_E\|_3^3
    &=3\int_0^1\lambda^2\sigma(S_\lambda)\,d\lambda \notag\\
    &\le4\pi N^{-3}
          +C_{p,\xi}N^{p-2+\xi}|E|
                                  \int_{1/N}^1\frac{d\lambda}{\lambda}
       \notag\\
    &\le C_{p,\xi}N^{p-2+\xi}(1+\log N)|E|.
    \label{eq:closure-indicator-strong}
\end{align}
The empty set is immediate.  This separately accounts for the entire
range \(\lambda<N^{-1}\).

\paragraph{Cell averages and general functions.}
Let \(f\in L^3(\R^3)\) be complex-valued.  The case \(f=0\) is
immediate, and homogeneity reduces the proof to \(\|f\|_3=1\).
Choose a dyadic \(N\ge2\) with
\(N^{-1}\le\delta<2N^{-1}\), and let \(F_N\) be the function equal
to the average of \(|f|\) on each cube of the \(N^{-1}\) grid.
Hölder's inequality on a cube and Jensen's inequality give
\begin{equation}
  0\le F_N\le N,\qquad \|F_N\|_3\le1.
  \label{eq:closure-cell-averages}
\end{equation}
Every grid cube meeting \(T_\delta(a,\omega)\) lies in the cylinder
with the same center and direction, length \(4\), and radius \(4/N\).
Indeed its points are within \(\sqrt3/N\) of a point of the original
cylinder, and \(N\ge2\).  Summing the integrals of \(|f|\) on all
such cubes therefore gives
\[
  \int_{T_\delta(a,\omega)}|f|
       \le\int_{\text{enlarged cylinder}}F_N.
\]
The enlarged volume is \(64\pi N^{-2}\), while
\(\pi\delta^2\ge\pi N^{-2}\).  Consequently
\begin{equation}
  K_\delta f\le64\mathcal K_N F_N.
  \label{eq:closure-enlargement}
\end{equation}

The part of \(F_N\) below \(N^{-1}\) contributes at most
\(N^{-1}\) pointwise to \(\mathcal K_N F_N\).  For its remaining
values use the disjoint level sets
\[
  E_j=\{2^j\le F_N<2^{j+1}\},
  \qquad -\log_2N\le j\le\log_2N.
\]
There are \(J_N=2\log_2N+1\) such levels.  Each is a finite union
of fine grid cubes: its volume is finite by
\eqref{eq:closure-cell-averages}, and every occupied cube has volume
\(N^{-3}\).  Sublinearity, the triangle inequality in \(L^3(S^2)\),
and \eqref{eq:closure-indicator-strong} yield
\begin{align*}
  \|\mathcal K_N F_N\|_3
    &\le (4\pi)^{1/3}N^{-1}
       +2C_{p,\xi}N^{(p-2+\xi)/3}(1+\log N)^{1/3}
                      \sum_j2^j|E_j|^{1/3}.
\end{align*}
Hölder's inequality for this finite sum gives
\[
  \sum_j2^j|E_j|^{1/3}
     \le J_N^{2/3}\left(\sum_j2^{3j}|E_j|\right)^{1/3}
     \le J_N^{2/3}\|F_N\|_3\le J_N^{2/3}.
\]
Together with \eqref{eq:closure-enlargement}, this proves
\[
  \|K_\delta f\|_3
     \le C_{p,\xi}N^{(p-2+\xi)/3}(1+\log N).
\]
Given \(\kappa>0\), take \(\xi=\kappa\) and absorb the logarithm
into \(C_\kappa N^{2\kappa/3}\).  Since \(N<2/\delta\), this is
\eqref{eq:closure-quantitative-maximal} for normalized \(f\).
Homogeneity gives the general case.  Fixed bounded resolutions are
included by enlarging the constant; directly, Hölder gives
\(K_\delta f\le(\pi\delta^2)^{-1/3}\|f\|_3\), uniformly in
direction, when \(\delta\) is bounded below.

\paragraph{Measurability.}
For clarity, all suprema used above define measurable functions of
direction.  If \((a_n,\omega_n)\to(a,\omega)\), the indicators of
the corresponding cylinders converge in \(L^{3/2}(\R^3)\): their
boundaries have measure zero and their supports lie in one bounded set
for a convergent sequence of parameters.  Hölder's inequality therefore
shows that the integral of \(|f|\) over the cylinder is continuous
in its center and direction.  For each fixed direction, its supremum
over all centers equals its supremum over the countable dense set
\(\mathbb Q^3\).  This is a countable supremum of continuous functions
of direction, hence is measurable.  The same reasoning applies to
\(\mathcal K_N\) and to the finite-volume indicator functions used
above.  It also shows that changing \(f\) on a null set changes none
of the tube integrals.  Thus all distribution and \(L^3\) estimates
are valid for the stated \(L^3\) equivalence classes.
\end{proof}

\begin{proof}[Proof of Theorem~\ref{thm:main}]
Fix \(\varepsilon>0\).  Proposition~\ref{prop:critical-vanishing}
allows a choice of \(p>2\) with \(h(p,0)=0\) and
\((p-2)/3<\varepsilon/2\).  Apply
Proposition~\ref{prop:maximal-reduction} with
\(\kappa=\varepsilon/2\).  Since \(0<\delta<1\),
\[
  \|K_\delta f\|_{L^3(S^2,\sigma)}
     \le C_{p,\varepsilon/2}
               \delta^{-\varepsilon}\|f\|_{L^3(\R^3)}.
\]
Choose this \(p\) once as a function of \(\varepsilon\) and rename
the constant \(C_\varepsilon\).  It is independent of \(\delta\),
of the complex-valued function \(f\), and of all tube centers and
directions, as required.
\end{proof}
